\subsection{Speech Self-Supervised Model Fine-Tuning}
Recent work has explored strategies to optimize speech SSL models for downstream tasks through architectural and training adaptations. \cite{zaiem2024speech} demonstrated that scaling probing head size not only improves model-specific performance on SER and SV tasks but also reshuffles SSL model rankings in benchmark evaluations. Work by \cite{zaiem2024less} investigated continual learning in SSL fine-tuning for ASR, revealing that full model unfreezing induces pre-training knowledge forgetting, whereas parameter-efficient methods like LoRA and Elastic Weight Consolidation mitigate this issue and achieve superior word error rates (WER), particularly on out-of-distribution data. Complementary to these approaches, \cite{peng2022ieee} proposed Multi-Head-Factorized Attentive Pooling (MHFA) for SV to hierarchically aggregate low-level speaker features with high-level phonetic information, thereby enhancing feature disentanglement.

\subsection{Layer Analysis of Speech Self-Supervised Models}
\cite{chen2022wavlm} found that layers at the beginning and middle of the WavLM model had higher magnitudes of learned weights after fine-tuning on speaker-dependent tasks, indicating their importance for speaker verification (SV) and speaker identification (SI). \cite{toyota23} used CCA analysis to show that models trained to predict discrete units learn phonetic and word information in intermediate layers, with this information concentrated in higher layers. They also demonstrated that layer-wise phone and word content correlate well with downstream task performance for phoneme recognition (PR) and automatic speech recognition (ASR). \cite{Yang24Foundations} showed that weights learned via weighted sum are not correlated with layers' performance on tasks. Both \cite{Yang24Foundations} and \cite{toyota23} found that for some tasks and SSL models, the best individual layer outperforms the learned weighted sum. However, this improvement in single-layer benchmarking is specific to individual SSL models and is not consistent across all models~\cite{toyota23}.

\subsection{Variational Inference and Early Exit}

From one perspective, our method could be seen as VAE \cite{vae_kingma} with layer index $l$ as a latent variable. Therefore, we can successfully adapt some techniques for it, like \cite{higgins2017betavae}. \cite{banino2021pondernet} used a similar approach, but for adaptive computational time \cite{graves2016adaptive} inference. This method was also adapted for early exiting during inference \cite{palbert}. While reducing inference time of the models is a promising direction, in our work, we concentrated solely on performance improvement.